\label{sec:model}

The Hilbert space Gaussian process (HSGP) construction of
\citet{solin2020hsgp} approximates a stationary GP by a finite
expansion in Laplacian eigenfunctions with Gaussian coefficients.
The domain determines the basis functions, while the covariance
kernel determines the coefficient variances. Changing the kernel or
its parameters therefore leaves the basis and its derivatives
unchanged, allowing their values at observation and prediction
locations to be computed once and reused
\citep{riutort2023practicalhsgp}. For shape-constrained regression,
a further advantage is that the response, slope, and curvature are
linear functions of the same coefficients. We can therefore express
all three restrictions in a common coefficient space and obtain
their conditioned laws from the ordinary Gaussian regression posterior.

Let $k(u)$ be a real stationary covariance function with a spectral
density, where $u=x-x'$ is the separation between inputs. We write
$S_k:\R\to[0,\infty)$ for its nonnegative, even, integrable spectral
density and use the Fourier convention
\[
 k(u)=\frac1{2\pi}\int_\R e^{\mathrm i\omega u}S_k(\omega)\,d\omega,
 \qquad u\in\R.
\]
Here $\omega$ is angular frequency, and $S_k(\omega)$ is the spectral
density associated with the covariance function $k$ through this
Fourier relation, evaluated at frequency $\omega$. The HSGP expansion
uses this density to assign variances at a discrete set of frequencies
determined by the Laplacian on a bounded interval.

Fix a center $c\in\R$, a half-width $L>0$, an integer $m\geq1$,
and a compact interval $I=[x_-,x_+]$ of positive length inside
$(c-L,c+L)$. On the padded interval $[c-L,c+L]$, the negative
Laplacian with zero boundary values has eigenvalues $\omega_j^2$
and orthonormal eigenfunctions $\phi_{j,L}$. Using the spectral
density $S_k$, these functions define the
finite Gaussian expansion
\begin{equation}\label{eq:finite-model}
 \begin{aligned}
 \phi_{j,L}(x)&=L^{-1/2}\sin(\omega_j(x-c+L)),
 &\omega_j&=j\pi/(2L),\\
 f_{\boldsymbol a}(x)&=a_0+\sum_{j=1}^ma_j\phi_{j,L}(x),
 &a_0&\sim N(\mu_0,\sigma_0^2),\quad a_j\sim N(0,S_k(\omega_j)).
 \end{aligned}
\end{equation}
The coefficients are independent, with $\sigma_0>0$ and
$S_k(\omega_j)>0$ at each retained frequency, so the finite coefficient
law is nondegenerate.


Under this prior, the covariance is
\[
 \Cov(f_{\boldsymbol a}(x),f_{\boldsymbol a}(x'))
 =\sigma_0^2+\sum_{j=1}^m S_k(\omega_j)
                         \phi_{j,L}(x)\phi_{j,L}(x').
\]
The sum approximates the stationary covariance $k(x-x')$, while
$\sigma_0^2$ is the contribution of the random intercept. Because
the eigenfunctions vanish at the padded endpoints, the finite
covariance is generally nonstationary. Defining the basis on an
interval larger than $I$ places this artificial boundary outside
the region of inference without requiring the response to vanish
at the endpoints of $I$. The intercept permits uncertainty in the
overall response level without changing slope or curvature.

\label{sec:finite-model}
For observations $\boldsymbol y\in\R^n$ at $x_i\in I$,
the Gaussian observation model is
\[
 \boldsymbol Y\mid\boldsymbol a\sim
 N\bigl((f_{\boldsymbol a}(x_1),\ldots,f_{\boldsymbol a}(x_n))^T,
        \boldsymbol\Sigma\bigr),
\]
where the noise covariance $\boldsymbol\Sigma$ is positive
definite and the errors are independent of the coefficients.
For given covariance and noise parameters and a chosen basis, we study
the coefficient posterior conditional on the observed responses and
their locations.


Write $\boldsymbol a=(a_0,\ldots,a_m)^T$. If $\boldsymbol X$ has rows
$(1,\phi_{1,L}(x_i),\ldots,\phi_{m,L}(x_i))$, set
$\boldsymbol m_0=(\mu_0,0,\ldots,0)^T$ and
$\boldsymbol V_0=\operatorname{diag}(\sigma_0^2,S_k(\omega_1),\ldots,S_k(\omega_m))$.
The coefficients and observations then have the joint Gaussian distribution
\[
 \begin{pmatrix}\boldsymbol a\\\boldsymbol Y\end{pmatrix}
 \sim N\left(
 \begin{pmatrix}\boldsymbol m_0\\\boldsymbol X\boldsymbol m_0\end{pmatrix},
 \begin{pmatrix}
  \boldsymbol V_0 & \boldsymbol V_0\boldsymbol X^T\\
  \boldsymbol X\boldsymbol V_0 &
  \boldsymbol X\boldsymbol V_0\boldsymbol X^T+\boldsymbol\Sigma
 \end{pmatrix}
 \right).
\]
Here $\boldsymbol\Sigma$ is the noise covariance, while
$\boldsymbol X\boldsymbol V_0\boldsymbol X^T+\boldsymbol\Sigma$
is the marginal covariance of the observations. Conditioning this joint
distribution on $\boldsymbol Y=\boldsymbol y$ gives the ordinary Gaussian posterior
\begin{equation}\label{eq:gaussian-posterior}
 \begin{aligned}
 \boldsymbol a\mid\boldsymbol y&\sim N(\boldsymbol m,\boldsymbol V),\\
 \boldsymbol m&=\boldsymbol m_0+\boldsymbol V_0\boldsymbol X^T
       (\boldsymbol X\boldsymbol V_0\boldsymbol X^T+\boldsymbol\Sigma)^{-1}
       (\boldsymbol y-\boldsymbol X\boldsymbol m_0),\\
 \boldsymbol V&=\boldsymbol V_0-\boldsymbol V_0\boldsymbol X^T
       (\boldsymbol X\boldsymbol V_0\boldsymbol X^T+\boldsymbol\Sigma)^{-1}
       \boldsymbol X\boldsymbol V_0.
 \end{aligned}
\end{equation}
These equations specify the Gaussian law that we
will condition on shape, with factorizations and linear solves used
in its numerical preparation.

We treat positivity, monotonicity, and convexity as sign restrictions
of orders zero, one, and two. Write $C^0(I)=C(I)$ and
$\frac{\mathrm{d}^0f}{\mathrm{d}x^0}=f$. For $r\in\{0,1,2\}$ we condition the ordinary
posterior on the function-space event
\begin{equation}\label{eq:shape-event}
 \mathcal C_r=\left\{f\in C^r(I):
                 \min_{x\in I}\frac{\mathrm{d}^rf}{\mathrm{d}x^r}(x)\geq0\right\}.
\end{equation}
The zero order gives a nonnegative response, the first gives
nondecreasing regression, and the second gives convex regression, with
$\mathcal C_r$ also denoting its explicitly identified preimage in
coefficient space. The restriction requires nonnegativity but does
not impose a fixed positive lower bound on the response, slope, or
curvature.

Shape information can equivalently be incorporated into the prior.
Provided that $\mathcal C_r$ has positive prior probability, we may
restrict the Gaussian prior to curves in $\mathcal C_r$ and then update
it using the observation likelihood. This gives the same posterior as
first performing ordinary Gaussian updating and then conditioning on
$\mathcal C_r$, because the prior normalizing constant cancels in
Bayes' formula.
